% 基本信息--------------------------------------------------------------------------------
% 指定论文类型0:毕业论文  1/2:专业实践1/2
\newcommand{\reportStyle}{0}

% 在这里填写你的论文中文题目
\newcommand{\thesisTitle}{基于区块链的去中心化音乐社区平台设计与实现}
% 在这里填写你的论文英文题目
\newcommand{\thesisTitleEN}{Design and Implementation of a Blockchain-Based Decentralized Music Community Platform}


% 若有副标题，则运行下一行的代码，若无副标题，则将下一行注释掉（在\haveSub{}最前面添加 % 号）
% \haveSub{}

% 在这里填写你的论文中文副标题（没有只需注释掉\haveSub{}即可）
\newcommand{\thesisSubTitle}{融合链上音乐资源检索与排序模块}
% 在这里填写你的论文英文副标题（没有只需注释掉\haveSub{}即可）
\newcommand{\thesisSubTitleEN}{With an On-chain Music Resource Retrieval and Ranking Module}

% 在这里填写你的相关信息
\newcommand{\deptName}{信息技术与人工智能学院}
\newcommand{\majorName}{软件工程}
\newcommand{\yourName}{陈普阳}
\newcommand{\yourStudentID}{220110790202}
\newcommand{\mentorName}{陈伟锋}
\newcommand{\className}{22软件2班}
\newcommand{\Today}{2026年5月}
% 基本信息--------------------------------------------------------------------------------

% 中英文摘要与关键词
\newcommand{\abstractCN}{
    随着数字音乐平台与区块链技术的发展，音乐作品在发布、确权、收益分配和访问授权等环节仍存在可信记录不足、流程透明性有限和实现成本较高等问题。本文设计并实现了基于区块链的链上音乐发布与播放系统 FreeFlow。系统采用”链上可信记录+链下高效服务”的 Web2.5 混合架构，使用智能合约管理作品发布、ERC-1155 访问凭证、购买授权与收益分账，结合 IPFS 保存音频、封面和元数据，并实现发布、检索、购买、播放、评论与链上音乐库等功能。为支持冷启动阶段的资源发现，本文实现了轻量检索模块。实验结果表明，4 项合约测试全部通过；在基于公开 MusicBrainz 样本库构建的 6000 条资源、162 个查询的评估集上，完整排序算法的 MRR 为 0.7253，nDCG@10 为 0.6480，零结果率为 3.70\%。该系统为区块链音乐平台的工程实践提供了一套可复现的原型参考，在音乐版权管理、创作者自主发行与收益透明分配等方向具有较广泛的应用前景。
}

\newcommand{\keywordsCN}{
    链上音乐；区块链；Web2.5；IPFS；访问控制；资源检索
}
\newcommand{\abstractEN}{
    With the development of digital music platforms and blockchain technology, music publishing, copyright verification, revenue distribution, and access authorization still face problems such as insufficient trustworthy records, limited process transparency, and high implementation cost. This thesis designs and implements FreeFlow, a blockchain-based on-chain music publishing and playback system. The system adopts a Web2.5 hybrid architecture featuring trustworthy on-chain records and efficient off-chain services. Smart contracts are used to manage music publishing, ERC-1155 access credentials, purchase authorization, and revenue splitting, while IPFS stores audio files, cover images, and metadata. The system supports publishing, retrieval, purchase, playback, comments, and on-chain music library management. To support resource discovery during the cold-start stage, a lightweight retrieval module is further implemented. Experimental results show that all four smart-contract tests pass. On an evaluation benchmark built from the public MusicBrainz sample database with 6000 resources and 162 queries, the complete ranking algorithm achieves an MRR of 0.7253, an nDCG@10 of 0.6480, and a zero-result rate of 3.70\%. The system provides a reproducible prototype reference for the engineering practice of blockchain-based music platforms, with broad application prospects in areas such as music copyright management, creator autonomy, and transparent revenue distribution.}
}

\newcommand{\keywordsEN}{
    On-chain music; Blockchain; Web2.5; IPFS; Access control; Resource retrieval
}
